
\documentclass{article}
\usepackage{colortbl}
\usepackage{makecell}
\usepackage{multirow}
\usepackage{supertabular}

\begin{document}

\newcounter{utterance}

\centering \large Interaction Transcript for game `adventuregame', experiment `home\_deliver\_three\_basic\_hard\_invlimittwo', episode 0 with reka{-}flash{-}3.1{-}t0.0.
\vspace{24pt}

{ \footnotesize  \setcounter{utterance}{1}
\setlength{\tabcolsep}{0pt}
\begin{supertabular}{c@{$\;$}|p{.15\linewidth}@{}p{.15\linewidth}p{.15\linewidth}p{.15\linewidth}p{.15\linewidth}p{.15\linewidth}}
    \# & \multicolumn{2}{c}{Player} && \multicolumn{2}{c}{Game Master} \\
    \hline

    \theutterance \stepcounter{utterance}  
    & & & \multicolumn{4}{p{0.6\linewidth}}{
        \cellcolor[rgb]{0.9,0.9,0.9}{
            \makecell[{{p{\linewidth}}}]{
                \texttt{\tiny{[P1$\langle$GM]}}
                \texttt{You are playing a text adventure game. I will describe what you can perceive in the game. You write the single action you want to take in the game starting with >. Only reply with an action.} \\
\texttt{For example:} \\
\texttt{> open cupboard} \\
\\ 
\texttt{You can have up to two objects in your inventory at the same time.} \\
\\ 
\texttt{Your goal for this game is: Put the pillow on the counter, the book on the shelf and the plate on the table.} \\
\\ 
\texttt{Once you have achieved your goal, write "> done" to end the game.} \\
\\ 
\texttt{You are in a broom closet now. There are a broom and a mop. The broom is on the floor. The mop is on the floor. There is a passage to a hallway here.} \\
            }
        }
    }
     \\ \\

    \theutterance \stepcounter{utterance}  
    & \multicolumn{4}{p{0.6\linewidth}}{
        \cellcolor[rgb]{0.561,0.737,0.561}{
            \makecell[{{p{\linewidth}}}]{
                \texttt{\tiny{[P1$\rangle$GM]}}
                \texttt{<reasoning>} \\
\texttt{I need to figure out how to get the pillow on the counter, the book on the shelf, and the plate on the table. Let's start by looking at the current situation. The player is in a broom closet with a broom and a mop on the floor. There's a passage to a hallway. The inventory can hold up to two objects.} \\
\\ 
\texttt{First, I need to move to the hallway to access other areas where the counter, shelf, and table might be. So the first action should be to go to the hallway. Let me check the example given: "open cupboard" is an action. So moving would be something like "> go to hallway".} \\
\\ 
\texttt{Once in the hallway, where else can I go? The problem mentions a counter, shelf, and table, which are likely in different rooms. Maybe the hallway connects to other rooms. Let's assume there's a kitchen with a counter, a living room with a shelf, and a dining area with a table. But I need to explore.} \\
\\ 
\texttt{Wait, the initial description only mentions the broom closet and the hallway. So maybe the hallway leads to other rooms. Let me think. The player needs to collect the pillow, book, and plate. Are these items in the broom closet? The closet has a broom and mop. So maybe the pillow, book, and plate are in other rooms connected via the hallway.} \\
\\ 
\texttt{So first step: go to the hallway. Then, explore the hallway to find other rooms. Let's say "> go to hallway". Then, perhaps there's a door or passage leading to the kitchen (counter), living room (shelf), and dining room (table). But how to navigate? Maybe the hallway has doors labeled or directions. Since it's a text adventure, the game might describe available exits.} \\
\\ 
\texttt{Assuming that in the hallway, there are passages to different rooms. Let's say "> go east" to the kitchen, "> go west" to the living room, etc. But without knowing the exact layout, I need to proceed step by step.} \\
\\ 
\texttt{Wait, the initial state is the broom closet. So first action is to go to the hallway. Let's do that. So first action: "> go to hallway".} \\
\\ 
\texttt{Once in the hallway, the game might describe other rooms. Let's say there's a kitchen to the north, a living room to the east, and a dining room to the west. The player needs to go to each room to} \\
            }
        }
    }
    & & \\ \\

    \theutterance \stepcounter{utterance}  
    & & & \multicolumn{2}{p{0.3\linewidth}}{
        \cellcolor[rgb]{0.9,0.9,0.9}{
            \makecell[{{p{\linewidth}}}]{
                \texttt{\tiny{[GM$|$GM]}}
                \texttt{Round: 0} \\
            }
        }
    }
    & & \\ \\

    \theutterance \stepcounter{utterance}  
    & & & \multicolumn{2}{p{0.3\linewidth}}{
        \cellcolor[rgb]{0.9,0.9,0.9}{
            \makecell[{{p{\linewidth}}}]{
                \texttt{\tiny{[GM$|$GM]}}
                \texttt{command\_tag\_missing} \\
            }
        }
    }
    & & \\ \\

    \theutterance \stepcounter{utterance}  
    & & & \multicolumn{2}{p{0.3\linewidth}}{
        \cellcolor[rgb]{0.9,0.9,0.9}{
            \makecell[{{p{\linewidth}}}]{
                \texttt{\tiny{[GM$|$GM]}}
                {'goal_states_achieved': [], 'game_successfully_finished': False}
            }
        }
    }
    & & \\ \\

    \theutterance \stepcounter{utterance}  
    & & & \multicolumn{2}{p{0.3\linewidth}}{
        \cellcolor[rgb]{0.9,0.9,0.9}{
            \makecell[{{p{\linewidth}}}]{
                \texttt{\tiny{[GM$|$GM]}}
                \texttt{command\_tag\_missing} \\
            }
        }
    }
    & & \\ \\

\end{supertabular}
}

\end{document}
